\begin{papercard}{wp292018}
\cf{Stages} \stagetag{S} (the legal meaning of transparency in the EU), \stagetag{C} (criteria for coding the clarity, accessibility and content of privacy notices), \stagetag{W} (where and when information must appear in an app), \stagetag{H} (duty to state consequences of processing). Article 29 Working Party, ``Guidelines on transparency under Regulation 2016/679'' (WP260 rev.01), adopted 29 November 2017, last revised 11 April 2018; 40 PDF pages \pc{wp292018}{1}{Guidelines on transparency under Regulation 2016/679}. Interpretative guidance, not a method or a study: it says what an EU regulator expects of controllers under GDPR Articles 12--14. Page 35 of the PDF (start of the Annex) has no text layer, so the annex rows are cited from pages 36--40.

\cf{Scope and purpose} Transparency is ``an overarching obligation'' covering (1) information to data subjects about fair processing, (2) communication about their rights and (3) facilitation of their rights \pc{wp292018}{4}{Transparency is an overarching obligation under the GDPR applying to three central areas}. It is about ``engendering trust in the processes which affect the citizen by enabling them to understand, and if necessary, challenge those processes'' \pc{wp292018}{4}{It is about engendering trust in the processes which affect the citizen by enabling them to understand, and if necessary, challenge those processes}. The concept ``is user-centric rather than legalistic'', and ``the quality, accessibility and comprehensibility of the information is as important as the actual content'' \pc{wp292018}{5}{the quality, accessibility and comprehensibility of the information is as important as the actual content of the transparency information}. The guidelines are generic and cannot address sector nuances. Transparency applies throughout the processing life cycle: at the start, throughout, and at specific points such as breaches or material changes \pc{wp292018}{6}{transparency applies at the following stages of the data processing cycle}. GDPR does not define transparency; Recital 39 gives the meaning \pc{wp292018}{6}{Transparency is not defined in the GDPR}.

\cf{Article 12 quality requirements} Information must be ``concise, transparent, intelligible and easily accessible'', in clear and plain language, in writing or other means (with oral provision on request) and free of charge \pc{wp292018}{6}{it must be concise, transparent, intelligible and easily accessible}.
\begin{itemize}\itemsep0pt
\item \emph{Concise and transparent}: present information efficiently ``to avoid information fatigue'' \pc{wp292018}{7}{present the information/ communication efficiently and succinctly in order to avoid information fatigue}, and keep it ``clearly differentiated'' from contractual provisions or general terms of use \pc{wp292018}{7}{This information should be clearly differentiated from other non-privacy related information such as contractual provisions or general terms of use}.
\item \emph{Intelligible}: understood ``by an average member of the intended audience''; controllers can test via user panels and readability testing \pc{wp292018}{7}{they can test these, for example, through mechanisms such as user panels, readability testing}. The data subject ``should not be taken by surprise at a later point'' and should be able to determine the scope and consequences in advance \pc{wp292018}{7}{the data subject should be able to determine in advance what the scope and consequences of the processing entails and that they should not be taken by surprise at a later point}. For complex, technical or unexpected processing controllers ``should also separately spell out in unambiguous language what the most important consequences of the processing will be'' \pc{wp292018}{7}{controllers should also separately spell out in unambiguous language what the most important consequences of the processing will be}.
\item \emph{Easily accessible}: the data subject ``should not have to seek out the information'' \pc{wp292018}{8}{The easily accessible element means that the data subject should not have to seek out the information}. For \textbf{apps}: information must be available ``from an online store prior to download'' and from inside the app, never more than ``two taps away'' \pc{wp292018}{8}{For apps, the necessary information should also be made available from an online store prior to download}, and ``specific to the particular app'' rather than the company's generic policy \pc{wp292018}{8}{the privacy information in question should be specific to the particular app and should not merely be the generic privacy policy of the company that owns the app}.
\item \emph{Clear and plain language}: information should be concrete and definitive, not ``abstract or ambivalent'' \pc{wp292018}{8}{it should not be phrased in abstract or ambivalent terms}. Poor examples: ``We may use your personal data to develop new services'', ``\ldots\ for research purposes'', ``\ldots\ to offer personalised services'' \pc{wp292018}{9}{We may use your personal data to develop new services}; good examples name the data type and what the processing does (for example keeping a record of clicked articles to target advertising on the site). Qualifiers such as ``may'', ``might'', ``some'', ``often'' and ``possible'' ``should also be avoided'' \pc{wp292018}{9}{Language qualifiers such as may might some often and possible should also be avoided}, and where used the controller must be able to show why they could not be avoided \pc{wp292018}{9}{they should be able, in accordance with the principle of accountability, to demonstrate why the use of such language could not be avoided}. Prefer the active voice \pc{wp292018}{10}{Writing should be in the active instead of the passive form}.
\item \emph{Children and vulnerable people}: language must suit children; transparency is ``a free-standing right'' that applies to children as to adults, and consent by a parent does not remove it \pc{wp292018}{10}{transparency is a free-standing right which applies as much to children as it does to adults}; vulnerabilities of other groups (such as people with disabilities) must be considered \pc{wp292018}{11}{if a data controller is aware that their goods/ services are availed of by (or targeted at) other vulnerable members of society}.
\end{itemize}

\cf{Form, layers and timing} Written form is the default, and a layered online notice is recommended, but the whole information should also exist ``in one single place or one complete document'' \pc{wp292018}{11}{the entirety of the information addressed to data subjects should also be available to them in one single place or one complete document}. Information may not depend on payment, and for free services the Article 13 information must come before sign-up \pc{wp292018}{13}{where free services are being provided to the data subject, the Article 13 information must be provided prior to, rather than after, sign-up}. The GDPR prescribes no format; the controller must choose ``appropriate measures'' in light of the device and user journey \pc{wp292018}{14}{The GDPR does not prescribe the format or modality by which such information should be provided to the data subject} and may test modalities with users \pc{wp292018}{14}{data controllers may wish to trial different modalities by way of user testing}. Article 13 covers data collected from the person, including by observation (cameras, Wi-Fi tracking, RFID, sensors); Article 14 data from other sources \pc{wp292018}{15}{a data controller collects from a data subject by observation}. For indirectly obtained data information must come within a reasonable period and at most one month \pc{wp292018}{16}{the maximum time limit within which Article 14 information must be provided to a data subject is one month}. A footnote requires that when a purpose includes creating \emph{inferred} personal data, the purpose and the categories of inferred data be communicated at collection \pc{wp292018}{14}{If the purpose includes the creation of inferred personal data, the intended purpose of creating and further processing such inferred personal data, as well as the categories of the inferred data processed, must always be communicated to the data subject}. Layered notices ``are not merely nested pages that require several clicks'' \pc{wp292018}{19}{layered privacy statements/ notices are not merely nested pages that require several clicks to get to the relevant information}; the \emph{first layer} should show purposes, controller identity and rights \pc{wp292018}{19}{the first layer/ modality should include the details of the purposes of processing, the identity of controller and a description of the data subject’s rights} and ``the processing which has the most impact on the data subject and processing which could surprise them'' \pc{wp292018}{19}{the first layer/ modality should also contain information on the processing which has the most impact on the data subject and processing which could surprise them}. Push (just-in-time) and pull (dashboards, permission management, tutorials) notices are named as user-centred options \pc{wp292018}{20}{Push notices involve the provision of just-in-time transparency information notices while pull notices facilitate access to information by methods such as permission management, privacy dashboards and learn more tutorials}.

\cf{Changes to the notice} The same principles apply to later changes. Always communicate a change of purpose, controller or rights procedure \pc{wp292018}{17}{Changes to a privacy statement/ notice that should always be communicated to data subjects include inter alia: a change in processing purpose; a change to the identity of the controller; or a change as to how data subjects can exercise their rights}; telling users to check the notice regularly for changes is ``insufficient but also unfair'' \pc{wp292018}{17}{should regularly check the privacy statement/notice for changes or updates are considered not only insufficient but also unfair}. Fundamental changes (more recipients, transfers to a third country) need notice well in advance \pc{wp292018}{17}{enlargement of the categories of recipients or introduction of transfers to a third country}, and transparency does not ``whitewash'' processing that has become different in nature \pc{wp292018}{18}{compliance with transparency requirements does not whitewash a situation where the changes to the processing are so significant}. Further processing for a new purpose needs notice beforehand, with a reasonable period that is longer ``the more intrusive (or less expected) the further processing'' \pc{wp292018}{24}{The principle of fairness requires that the more intrusive (or less expected) the further processing, the longer the period should be}.

\cf{Tension, profiling, risks, icons, rights} The text names ``an inherent tension'' between comprehensiveness and conciseness \pc{wp292018}{18}{There is an inherent tension in the GDPR between the requirements on the one hand to provide the comprehensive information to data subjects}. Profiling and automated decisions require ``meaningful information about the logic involved and the significant and envisaged consequences'' \pc{wp292018}{22}{meaningful information about the logic involved and the significant and envisaged consequences of the processing for the data subject}; publishing a DPIA is suggested as a trust measure but is not obligatory \pc{wp292018}{22}{data controllers may consider publication of the DPIA (or part of it)}. Icons may accompany but not replace the required information \pc{wp292018}{25}{the use of icons should not simply replace information necessary for the exercise of a data subject’s rights}, and standard icons need prior research on efficacy \pc{wp292018}{26}{it will be necessary for extensive research to be conducted in conjunction with industry and the wider public as to the efficacy of icons in this context}. Rights: a triple obligation to inform about rights, to communicate clearly about them and to facilitate them \pc{wp292018}{26}{Transparency places a triple obligation upon data controllers insofar as the rights of data subjects under the GDPR are concerned}. Exemptions in Article 14.5 (impossible or disproportionate effort, legal obligation, professional secrecy) require a documented balancing of effort against impact on the person \pc{wp292018}{31}{carry out a balancing exercise to assess the effort involved for the data controller to provide the information to the data subject against the impact and effects on the data subject}; a breach notification must meet the same clarity requirements \pc{wp292018}{34}{The communication of a data breach must satisfy the same requirements}.

\cf{Annex: what must be stated} The annex table summarises Articles 13 and 14. Legitimate interests must be identified specifically, with the balancing-test result available on request \pc{wp292018}{36}{The specific interest in question must be identified for the benefit of the data subject}. A ``recipient'' need not be a third party \pc{wp292018}{37}{a recipient does not have to be a third party}; naming recipients is generally the most meaningful \pc{wp292018}{37}{In practice, this will generally be the named recipients, so that data subjects know exactly who has their personal data}, otherwise categories ``as specific as possible'' by activity, industry, sector and location \pc{wp292018}{37}{the information should be as specific as possible by indicating the type of recipient (i.e. by reference to the activities it carries out), the industry, sector and sub-sector and the location of the recipients}. Transfers to third countries should name the country \pc{wp292018}{38}{this will generally mean that the third countries be named}. Retention: ``It is not sufficient \ldots\ to generically state that personal data will be kept as long as necessary'' \pc{wp292018}{38}{It is not sufficient for the data controller to generically state that personal data will be kept as long as necessary for the legitimate purposes of the processing}. The right to object must be brought to attention at the latest at first communication ``clearly and separately from any other information'' \pc{wp292018}{39}{the right to object to processing must be explicitly brought to the data subject’s attention at the latest at the time of first communication with the data subject and must be presented clearly and separately from any other information}, and withdrawal of consent must be ``as easy \ldots\ as to give it'' \pc{wp292018}{39}{it should be as easy for a data subject to withdraw consent as to give it}. Further rows: right to complain, whether providing data is statutory or contractual with consequences of refusal (online forms should mark required fields), the source of the data, and automated decision-making including profiling with its logic and consequences.

\cf{Limits (card)} Interpretative guidance of an advisory body, not case law or a measurement standard; sector-neutral; no evidence in the file on whether the recommended practices improve understanding (it calls for research on icons and recommends user testing); health data and consent under the special-category rules (GDPR Art.~9) are not treated here; the document addresses controllers, so ``good'' and ``poor'' examples are illustrations, not a coding scheme; the status of the guidelines after 2018 (adoption by the European Data Protection Board) is not in the file and should be checked before the thesis states it.

\cf{Use for this thesis} (1) Normative criteria for the ``communication quality'' codes: concise and separated from terms of use; intelligible for the intended audience; accessible before download and within two taps; app-specific; plain language; vague qualifiers (``may'', ``might'', ``some'', ``often'', ``possible'') and the quoted poor examples (``develop new services'', ``research purposes'', ``personalised services'') as seeds for a ``vague purpose'' code, with the good examples as contrast. (2) Paragraph 10 and the first-layer rule (para.~36) are the most direct regulatory counterpart to the thesis' question: they require stating the \emph{consequences} of complex or unexpected processing and the processing with most impact; a policy that lists categories but no consequences can be coded as not meeting that expectation, while noting that legal compliance is not the thesis' measure. (3) The footnote on inferred data and the rules on profiling support a code for inference-related disclosures. (4) The annex is a ready completeness checklist: purposes and legal basis, legitimate interests, named or specific recipients, third-country transfers and safeguards, retention not stated as ``as long as necessary'', rights, withdrawal of consent, complaint, statutory or contractual requirement, source, automated decisions, change notification. (5) The walkthrough (stage W) can test the store-before-download, in-app access (two taps), just-in-time notices and dashboards named here. (6) For the German DiGA context cite it with the BfArM guide and note the EU setting; compare with US-based studies in the corpus that assume different law.
\end{papercard}
